% SPDX-License-Identifier: MIT
\chapter{Troubleshooting}
\label{ch:troubleshooting}

\section{Installation}

\begin{description}[style=nextline,leftmargin=0pt,labelindent=0pt,labelwidth=0pt]
  \item[\cmd{Could not load the Qt platform plugin "xcb"} (Linux)]
    Missing Qt system libraries. Install them
    (Section~\ref{sec:linux-libs}), or use the conda environment, which supplies
    them as packages. To see exactly which library is missing, run
    \cmd{QT\_DEBUG\_PLUGINS=1 makeyourtree-studio}.
  \item[\cmd{OSError: [WinError 206] The filename or extension is too long}]
    The clone is too deep for Windows' \cmd{MAX\_PATH} limit. Delete
    \cmd{.venv}, re-clone somewhere short such as
    \cmd{C:\textbackslash dev\textbackslash makeyourtree}, and install again.
  \item[\cmd{ModuleNotFoundError: No module named 'makeyourtree'}, but \cmd{pytest} passes]
    The same \cmd{MAX\_PATH} failure, which rolls back partially. \cmd{pytest}
    still passes because its configuration adds \cmd{src/} to the path, masking
    the missing install.
  \item[\cmd{cannot be loaded because running scripts is disabled} (Windows)]
    PowerShell execution policy. Run
    \cmd{Set-ExecutionPolicy -Scope Process -ExecutionPolicy Bypass}, then
    activate again.
  \item[\cmd{ensurepip is not available} (Linux)]
    Install \cmd{python3.12-venv} on Debian and Ubuntu.
  \item[The environment will not solve; \cmd{pyside6} not found]
    You are on the \cmd{defaults} channel. Prefer conda-forge
    (Chapter~\ref{ch:installation}).
  \item[Qt appears twice, or the environment is enormous]
    A \cmd{pip install -e .} without \opt{--no-deps} inside a conda environment.
    Remove the environment and redo the install step.
\end{description}

\section{Files and parsing}

\begin{description}[style=nextline,leftmargin=0pt,labelindent=0pt,labelwidth=0pt]
  \item[``My file has problems'' but it looks fine]
    Most diagnostics are informational. The commonest,
    \cmd{newick.internal-label-ambiguous}, simply records that numeric internal
    labels were read as support values --- which is nearly always right. Open
    \keys{Ctrl+D} and read them; a clean file often produces one or two.
  \item[Support values appear as node names, or vice versa]
    The same ambiguity, resolved the other way. Numeric internal labels are read
    as support by default; pass \cmd{internal\_labels='name'} through the API if
    yours are genuinely names.
  \item[Annotation rows do not match any tip]
    Names must match the tree exactly. The import wizard shows a match report
    before committing --- read it. Common causes are underscores against spaces,
    trailing whitespace, and accession prefixes present in one file but not the
    other.
  \item[A tip shows a zero where it should show nothing]
    Check the source data: an empty field or \cmd{-} means \emph{no value} and
    draws a gap. If you see a bar of zero height instead, the file really does
    contain a zero.
\end{description}

\section{Figures}

\begin{description}[style=nextline,leftmargin=0pt,labelindent=0pt,labelwidth=0pt]
  \item[All the text is empty boxes]
    Qt found no fonts. This happens under the \cmd{offscreen} platform, which
    reports zero font families. Clear \cmd{QT\_QPA\_PLATFORM}, or use
    \cmd{xvfb-run} on a headless machine. Vector SVG from the core is unaffected
    because it stores text as text.
  \item[There is a gap between the tips and the annotation ring]
    In a circular \emph{phylogram} this is correct: tips genuinely sit at
    different radii. Turn on \textbf{Align tips} (\keys{A}) to close it.
  \item[Tracks in an unrooted layout look wrong]
    They are. An unrooted layout has no tip ordering for band space to use, so
    tracks are drawn as an unaligned column. Use a rooted layout
    (Section~\ref{sec:unrooted-tracks}).
  \item[The heatmap ignores my \opt{palette} setting]
    The heatmap has no \opt{palette} option; its ramp is set by
    \opt{color\_min}, \opt{color\_mid} and \opt{color\_max}. Unknown options are
    recorded as diagnostics rather than applied.
  \item[Labels are clipped at the edge of the figure]
    Increase \opt{width}, or limit label length with \opt{max\_label\_width}.
  \item[The figure is unreadably tall]
    Reduce \opt{row\_spacing}, turn off tip labels, or collapse clades
    (\keys{C}).
\end{description}

\section{Performance}

\begin{description}[style=nextline,leftmargin=0pt,labelindent=0pt,labelwidth=0pt]
  \item[The window freezes when I change layout on a big tree]
    Known limitation. Above a few thousand tips, switching layout takes several
    seconds during which the window does not repaint. It is working, not hung,
    and the operation is not currently cancellable.
  \item[Rendering a very large tree produces an enormous SVG]
    Expected: a 100{,}000-tip tree is about 28\,MB of SVG. Export PNG at a
    chosen scale instead if the file size matters more than resolution
    independence.
\end{description}

\section{Reporting a problem}

Please open an issue with your platform, your Python version, the full error
output, and --- if you can share it --- a small tree that reproduces the problem.
A synthetic tree that shows the same failure is ideal, since it avoids sharing
unpublished data.
